% 第29章思考题29.1--29.13独立解析；本文件不属于原书章末习题705对。
\section*{第29章\quad 固体中的电子：思考题解析}
\paragraph{思考题 29.1}
自由电子近似把金属价电子视为
在体积 \(V\) 内运动、服从费米统计且相互作用弱的
粒子。周期晶格势真实存在，
Bloch态 \(\psi_{\boldsymbol k}(\boldsymbol r)
=u_{\boldsymbol k}(\boldsymbol r)e^{i\boldsymbol k\cdot\boldsymbol r}\)
替代真正的自由平面波。
在费米面附近能带近似抛物
\(E(\boldsymbol k)\simeq E_0+
\hbar^2|\boldsymbol k|^2/(2m^*)\) 且散射较弱时，
才可用有效质量的准自由电子模型。
\[
\boxed{\text{“自由”指有效低能运动近似，
不是完全没有晶格作用。}}
\]

\paragraph{思考题 29.2}
室温 \(T\ll T_F\)，费米面深处的电子因
Pauli原理难以吸收能量，
只有费米面附近宽度 \(\sim k_BT\) 的
小部分电子参与热激发，故
\(C_{V,e}\propto T/T_F\)，远小于经典值。
但施加微弱电场时，费米面附近的Bloch电子
便能产生小的分布位移和净电流；
在Drude近似下
\[
\boxed{\sigma=ne^2\tau/m^*}.
\]
大量载流子和长弛豫时间可使
\(\sigma\) 很高。
\emph{电子总数多}不等于电子热容大。

\paragraph{思考题 29.3}
电子自旋1/2是费米子，Pauli原理要求
单量子态平均占据不超过1，
平衡分布为
\[
\boxed{f(E)=\frac1{\exp[(E-\mu)/(k_BT)]+1}}.
\]
在 \(T\to0\) 时成为阶跃
\(f(E)=1\ (E<E_F)\)、\(0\ (E>E_F)\)。
经典Maxwell--Boltzmann分布
\(f\propto e^{-(E-\mu)/(k_BT)}\)
仅在极低占据、非简并条件成立，
不可描述 \(T\ll T_F\) 金属。

\paragraph{思考题 29.4}
晶体中大量等价原子相互作用，
原先局域分立能级裂变成大量密集的准连续能态，
组成\emph{允许能带}；
相邻允许带之间不存在Bloch本征态的能量区间
称\emph{禁带}。
在半导体/绝缘体通常称最高满带为价带，
最低空的低能带为导带，
\[
\boxed{E_g=E_{C,\min}-E_{V,\max}}.
\]
但金属里有部分占据的带穿过费米能级，
因此不能普遍把“价带满、导带空”推广到所有固体。

\paragraph{思考题 29.5}
金属通常费米能级穿过一条允许带，
或两带重叠，因此有许多低能电荷激发。
绝缘体和半导体在 \(T=0\) 时
价带满、导带空，费米能级处于禁带附近；
两者区别主要是可热激发/光激发的
能隙大小与相应室温载流子密度。
掺杂半导体的化学势会向施主/受主能级移动，
不能仅用一个固定“费米能级在带隙正中央”
描述所有掺杂状态。

\paragraph{思考题 29.6}
四价Si晶格掺五价P，产生施主态；
温度允许施主电离时，
可动电子增加，成为\emph{n型半导体}。
对应正电的施主离子留在晶格中：
\[
\boxed{n+N_A^-=p+N_D^+}
\]
是一般电中性条件（\(N_D^+\) 为电离施主，
\(N_A^-\) 为电离受主）。
大量电子进入导带并不表示宏观晶体带负电；
在非简并、完全电离的施主主导情形
\(n\simeq N_D\)，少数空穴
\(p\simeq n_i^2/n\) 反而减少。

\paragraph{思考题 29.7}
三价硼掺四价Ge产生\emph{受主}，
受主俘获价带电子，使可动空穴成为多数载流子。
在非简并热平衡本征关系
\[
\boxed{np=n_i^2}
\]
下，若 \(p\) 增加，少数载流子导带电子
\(n=n_i^2/p\) 通常\emph{减少}，
并非同步增加。
俘获价电子的受主离子带负电，
与正空穴电荷补偿，宏观保持近似电中性，
不能把“空穴带正电”等同样品总体带正电。

\paragraph{思考题 29.8}
均匀本征或单一掺杂的半导体体材料，
微弱外电场下正反两向通常都有电流，
不是理想单向导体。
\emph{PN结}具有内建电势障碍：
正向偏压降低势垒、增加载流子注入，
反向偏压增加势垒而仅有较小漏电流，
因而表现出\emph{整流}。
它的整流性质来自\emph{非均匀空间结和势垒}，
不是“所有半导体材料天然只能单向导电”。

\paragraph{思考题 29.9}
单一载流子Drude模型的霍尔系数
\(R_H\simeq1/(nq)\)，
给定电流密度和磁场，
\[
\boxed{E_H=R_H jB},\qquad
V_H\simeq\frac{IB}{nqt}
\]
（薄片厚度 \(t\)，符号按载流电荷正负决定）。
通常金属载流子浓度 \(n\sim10^{29}\,\mathrm{m^{-3}}\)，
半导体的可调掺杂载流子浓度可小得多，
因而相同 \(I,B,t\) 下
\(|R_H|\)、\(|V_H|\) 大得多。
半导体中若电子与空穴共存需用双载流子公式。

\paragraph{思考题 29.10}
取东为 \(+\hat x\)、北为 \(+\hat y\)、
上为 \(+\hat z\)。题给正电流从北向南，
\(\boldsymbol j=-j\hat y\)，
\(\boldsymbol B=B\hat z\)。
N型主要载流子电子的漂移速度
\(\boldsymbol v_e\parallel+\hat y\)。
霍尔稳态时
\(\boldsymbol E_H=-\boldsymbol v\times\boldsymbol B\)，
因此
\[
\boxed{\boldsymbol E_H^{(n)}\parallel-\hat x
\quad\text{即向西}}.
\]
P型主要空穴
\(\boldsymbol v_h\parallel-\hat y\)，
得到
\[
\boxed{\boldsymbol E_H^{(p)}\parallel+\hat x
\quad\text{即向东}}.
\]
若磁场或\emph{电流方向}反转，两个结果都反向。
方向判断不能只依据“电子电荷负”而忽略
电子漂移方向与常规电流相反。

\paragraph{思考题 29.11}
本征半导体同时有电子与空穴，
霍尔响应并非必然消失。
在弱磁场、双载流子Drude近似，
\[
\boxed{
R_H=\frac{p\mu_h^2-n\mu_e^2}
{e(p\mu_h+n\mu_e)^2}}.
\]
本征时 \(n=p\)，故
\[
R_H=\frac{\mu_h^2-\mu_e^2}
{en(\mu_h+\mu_e)^2}.
\]
若\emph{迁移率相同}
\(\mu_h=\mu_e\)，两贡献正好抵消；
实际多数半导体的电子/空穴迁移率不同，
仍可测得非零霍尔电压，
其符号反映迁移率加权竞争。

\paragraph{思考题 29.12}
在原图29.17的N型沟道、沿源到漏
沟道电势逐渐升高的常见偏置下，
局部栅与沟道电位差
\[
V_{GC}(x)=V_G-V(x)
\]
沿靠近漏端变得更负。
提高源极相对栅极电压 \(U_{SG}\)
使耗尽区加宽时，局部可导电层通常
\emph{先在漏极侧达到耗尽}：
\[
\boxed{\text{先从漏端阻断，再沿沟道向源端发展}}.
\]
严格位置取决于器件是耗尽型还是增强型、
\(V_{DS}\) 的极性与大小及图29.17的定义；
当 \(V_{DS}=0\) 且栅偏压空间均匀，
可近似同时耗尽。不能未经检查原图的
\(U_{SG},I_{DS}\) 正负号套用所有MOSFET。

\paragraph{思考题 29.13}
遥控器发射端使用红外发光二极管（LED）
及相应的晶体管驱动，
通过编码调制发射红外光脉冲。
电视机接收端使用硅光电二极管
或光电晶体管、跨阻/放大器和解调电路，
把红外光子转换为电信号。
\[
\boxed{\text{发射：红外LED；接收：光电二极管/光电晶体管}}.
\]
现代器件封装可能把滤光、
放大和解调集成在单个接收模块中；
这不同于PN结本身“只能单向导电”的简单叙述。
